\section{理论框架}
\label{sec:framework}

在本节中，我们回顾本次计算所需的必要理论背景。我们给出
CS 核与快度演化的定义，并引入准 TMD 波函数。随后我们讨论准 TMDWFs 的因子化
及其与 CS 核的联系。

\subsection{Collins-Soper 核与快度演化}


与共线的光锥 PDFs 及分布振幅不同，TMDPDFs 与 TMDWFs 既依赖重正化标度 $\mu$，也依赖一个额外的快度重正化标度。后者之所以出现，是因为矩阵元还会遭受所谓的快度发散，需要专门的调节子~\cite{Collins:1981uk,Becher:2010tm,Chiu:2011qc}。
在 TMD 因子化中，硬模式（即远离壳的模式）对树图过程的贡献通常在维数正规化方案下计算。与沿不同方向高速运动的部分子相联系的共线模式，以及典型动量为 $\Lambda_{\mathrm{QCD}}$ 量级的软模式，具有相同的虚度，
仅能通过其快度加以区分。在使用维数正规化这类只调节紫外发散的方案进行计算时，对快度积分后会在软矩阵元与共线矩阵元中遇到额外的快度发散，必须用专门的调节子来处理。经过这一正规化之后，TMDPDFs 与 TMDWFs 获得额外的快度标度依赖。这种依赖在物理观测量的理论预言中应当抵消。

CS 核 $K\left(b_{\perp}, \mu\right)$ 即熟知的快度反常量纲，它编码了 TMD 分布的快度依赖~\cite{Collins:1981va,Collins:1981uk}：
\begin{align}
	2 \zeta \frac{d}{d \zeta} \ln f^{\mathrm{TMD}}\left(x, b_{\perp}, \mu, \zeta\right)=K\left(b_{\perp}, \mu\right), \label{eq:TMDsrapidityevolution}
\end{align}
其中 $f^{\mathrm{TMD}}$ 表示任意领头扭度的 TMDPDF 或 TMDWF。TMD 分布依赖于纵向动量分数 $x$、横向空间间隔 $b_{\perp}$（即横动量 $k_{\perp}$ 的傅里叶共轭量），以及重正化标度 $\mu$ 和快度标度 $\zeta${后者与强子动量相关}。CS 核 $K\left(b_{\perp}, \mu\right)$ 的 $\mu$ 依赖满足如下的重正化群方程（RGE）：
\begin{align}
	\mu^{2} \frac{d}{d \mu^{2}} K\left(b_{\perp}, \mu\right)=-\Gamma_{\text {cusp}}\left(\alpha_{s}\right).  \label{eq:cuspanomalousdimension}
\end{align}
这里 $\Gamma_{\text {cusp}}\left(\alpha_{s}\right)=\alpha_sC_F/\pi+\mathcal{O}(\alpha_s^2)$ 是尖点反常量纲，其在微扰论中已算到两圈（文献~\cite{Li:2016ctv}）与三圈（文献~\cite{Moch:2017uml}）。该 RGE 的解可表示为：
\begin{align}
	K\left(b_{\perp}, \mu\right)=-2 \int_{1 / b_{\perp}}^{\mu} \frac{\mathrm{d} \mu^{\prime}}{\mu^{\prime}} \Gamma_{\text {cusp}}\left(\alpha_{s}\left(\mu^{\prime}\right)\right)+K\left(\alpha_{s}\left(1 / b_{\perp}\right)\right).
	\label{eq:nonper-per-CS-kernel}
\end{align}
对于满足 $b_{\perp}^{-1}\lesssim\Lambda_{\mathrm{QCD}}$ 的大 $b_{\perp}$，CS 核变为非微扰的，由上式中的非尖点反常量纲 $K\left(\alpha_{s}\left(1/b_{\perp}\right)\right)$ 表示。

在过去几十年中，人们已在 TMD 部分子分布的全局拟合中对 CS 核做了广泛研究~\cite{Landry:1999an,Landry:2002ix,DAlesio:2014mrz,Sun:2014dqm,Konychev:2005iy,Bacchetta:2017gcc,Scimemi:2017etj,Scimemi:2019cmh,Bacchetta:2019sam}。非微扰区域的显式形式只能通过把小 $b_{\perp}$ 处的微扰表达式外推参数化得到，这不可避免地引入系统不确定度。直到 LaMET~\cite{Ji:2013dva,Ji:2014gla} 建立之前，在格点上直接计算 TMDPDFs 及相关的 CS 核几乎是一道不可逾越的障碍。LaMET 近期的一项显著进展是：这些量可以通过相应的可观测量来获取~\cite{Ebert:2019okf,Ji:2019sxk,Ji:2019ewn,Ji:2021znw}。

%%%%%%%%%%%%%%%%%%%%%%%%%%%%%%%


\subsection{准 TMD 波函数}
如上所述，可以为一个沿 $z$ 方向以大动量 $P^z$ 高速运动的赝标介子定义准 TMDWFs：
\begin{align}
	\tilde{\Psi}^{\pm}&\left(x,b_{\perp},\mu,\zeta_z\right)=\lim_{L\to\infty}\int\frac{dz}{2\pi}e^{ix_rzP^z}\frac{\tilde{\Phi}^{\pm0}\left(z,b_{\perp},P^z,a,L\right)}{\sqrt{Z_E(2L,b_{\perp},\mu, a)}},\label{eq:quasiWFinmomentumspace}
\end{align}
其中 $x_r=x-\frac{1}{2}$。未减除的准 TMDWF $\tilde{\Phi}^{\pm0}$ 定义为一个等时关联函数，它包含带有钉扎形（staple 形）规范链的非定域夸克双线性算符：
\begin{align}
	\tilde{\Phi}^{\pm0}&\left(z,b_{\perp},P^z,a,L\right) =\left\langle 0\right|\bar{\psi}\left(z \hat{n}_{z}/2+b_{\perp} \hat{n}_{\perp}\right) \Gamma \nonumber\\
	&\times U_{\sqsupset, \pm L}\left(z \hat{n}_{z}/2+b_{\perp} \hat{n}_{\perp},-z \hat{n}_{z}/2\right) \psi\left(-z \hat{n}_{z}/2\right)\left| P^{z}\right\rangle. \label{eq:quasiWFincoordinatespace}
\end{align}
对于赝标介子态，狄拉克结构 $\Gamma$ 可取为 $\gamma^z\gamma_5$ 或 $\gamma^t\gamma_5$，它们在光锥极限下趋于领头扭度结构 $\gamma^+\gamma_5$。当 $P^z$ 大而有限时，$\gamma^z\gamma_5$ 与 $\gamma^t\gamma_5$ 之差被按 $1/P^z$ 的幂次压低。在技术上，也可以采用两者的组合，例如 $\left(\gamma^z\gamma_5+\gamma^t\gamma_5\right)/2$，以使幂次修正最小化；更多细节见第~\ref{sec:operator}节。文献~\cite{Li:2021wvl}也探讨了多种组合。$\tilde{\Phi}^{\pm0}$ 中的上标 ''0'' 表示裸量。线性发散来自规范链的自能，
\begin{align}
	&U_{\sqsupset, \pm L}\left(z \hat{n}_{z}/2+b_{\perp} \hat{n}_{\perp},-z \hat{n}_{z}/2\right) \equiv\nonumber\\
	&	U_z^{\dagger}\left(z \hat{n}_{z}/2+b_{\perp} \hat{n}_{\perp}; L \right) U_{\perp}\big((L-z /2)\hat{n}_{z};{b}_T \big) U_z\left( -z \hat{n}_{z}/2; L \right), \label{eq:stapleshapedUlink}
\end{align}
且不以 $d=4$ 处极点的形式出现在维数正规化中。$U_{\sqsupset, \pm L}$ 中的欧几里得规范链定义为
\begin{align}
	U_z(\xi,\pm L)=\mathcal{P}	\exp \left[-i g \int_{\xi^{z}}^{\pm L} d \lambda n_{z} \cdot A\left(\vec{\xi}_{\perp}+n_{z} \lambda\right)\right],
\end{align}
其中 $\xi^{z}=-\xi \cdot n_{z}$。$\pm L$ 对应在有限的欧几里得格点上规范链沿正或负 $n_z$ 方向所能到达的最远位置。这在图~\ref{fig:definitionofquasiTMDWF}中以蓝线与红线描绘。


%%%%%%%%%%%%%%%%%%%%%%%%%%%%%%%%%%
\begin{figure}
\centering
\includegraphics[width=0.5\textwidth]{quasi-TMDWF}
\caption{未减除准 TMDWF 中所含钉扎形规范链及相关 Wilson 圈的示意。上图中蓝色与红色双线表示欧几里得格点上的 $L$ 平移方向，下图给出相应的 Wilson 圈，它将减除准 TMDWF 中的紫外对数发散与线性发散。}
\label{fig:definitionofquasiTMDWF}
\end{figure}
%%%%%%%%%%%%%%%%%%%%%%%%%%%%%%%%%%

由于线性发散与规范链相联系，可以用一条总长度相同的类似规范链将其去除。一个可选的做法是利用记为 $Z_E$ 的 Wilson 圈。Wilson 圈可以取为 $z$-$\perp$ 平面内平坦矩形欧几里得 Wilson 圈的真空期望值：
\begin{align}\label{eq:WilsonLoop}
Z_E\left(2L, b_{\perp}, \mu,a\right)=\frac{1}{N_{c}} \operatorname{Tr}\left\langle 0\left|U_{\perp}(0;b_{\perp}) U_{z}\left(b_{\perp}\hat{n}_{\perp}; 2L\right)\right| 0\right\rangle.
\end{align}
这里 $Z_E$ 在 $z$ 方向的长度是钉扎形规范链 $U_z^{(\dagger)}$ 的两倍，因此可以预期 $Z_E\left(2L, b_{\perp}, \mu\right)$ 的平方根会消去规范链中的线性发散与重夸克势。Wilson 线与轻夸克的顶点处存在残余的对数发散，可以在维数正规化中重正化~\cite{Ji:2021uvr}。由于这些对数发散与 $z$、$b_{\perp}$、$P_{z}$ 及 $L$ 无关，当研究准 TMDWFs 之比时它们将明显地相互抵消。



%%%%%%%%%%%%%%%%%%%%%%%%%%%%%%%%%%%%%%%%
\subsection{准 TMDWFs 的因子化}\label{sec:factorizationformula}

借助软函数，减除后准 TMDWFs 中的红外贡献可以被恰当地计入，使得准 TMDWFs 与光锥波函数的红外结构相互匹配。这意味着 LaMET 框架下一个乘法型的因子化定理~\cite{Ebert:2019okf,Ji:2019sxk,Ji:2019ewn,Ji:2021znw,Ebert:2022fmh}：
\begin{align}
&\tilde \Psi^{\pm}(x, b_{\perp}, \mu, \zeta_z) S_r^{1/2}(b_{\perp}, \mu) \nonumber\\
&= H^{\pm}\left(\zeta_z, \overline\zeta_z, \mu^2\right) \exp\left[\frac{1}{2}K(b_{\perp},\mu)\ln\frac{\mp\zeta_z- i\epsilon}{\zeta} \right]  \nonumber\\
& \qquad\times \Psi^{\pm}(x, b_\perp, \mu, \zeta)+\mathcal{O}\left(\frac{\Lambda_{QCD}^2}{\zeta_z}, M^2\zeta_z,\frac{1}{b_{\perp}^2\zeta_z}\right),
\label{eq:matching}
\end{align}
其中式~\eqref{eq:matching}中的上标 $\pm$ 对应 Wilson 线的方向，$\Psi^{\pm}$ 是在无限动量系中定义的 TMDWFs。约化软函数 $S^{1/2}_r(b_{\perp}, \mu)$ 源于 $\tilde \Psi^{\pm}$ 与 $\Psi^{\pm}$ 中不同的软胶子辐射效应~\cite{Ji:2021znw}。快度标度 $\zeta$ 与 $\zeta_z$ 的失配可以由 CS 核 $K(b_{\perp},\mu)$ 补偿。$S$ 与 $K$ 都不依赖 $\pm$ 的选择。$H^{\pm}$ 是一圈微扰匹配 kernel~\cite{Ji:2021znw}：
\begin{align}
& H^{\pm}(\zeta_z, \overline{\zeta}_z,\mu) \nonumber\\
& = 1+  \frac{ \alpha_s C_F}{4\pi} \bigg(-\frac{5\pi^2}{6}-4 +{\ell}_{\pm} + \overline \ell_{\pm} - \frac{1}{2} ({\ell}_{\pm}^2 + \overline \ell_{\pm}^2) \bigg).  \label{eq:1loophardkernel}
\end{align}
采用缩写 ${\ell}_{\pm} = \ln \left[(-\zeta_z \pm i\epsilon)/\mu^2\right]$、$\overline {\ell}_{\pm} = \ln \left[(-\overline \zeta_z \pm i\epsilon)/\mu^2\right]$，以及标度 $\zeta_z = (2x P^z)^2$、$\overline\zeta_z = \left(2\bar{x} P^z\right)^2$、$\bar{x}=1-x$。应当注意，$H^{\pm}$ 在 $\ell_{\pm}$ 与 $\bar{\ell}_{\pm}$ 中含有非零的虚部。虽然 $\ell_{\pm}/\bar{\ell}_{\pm}$ 中的虚部是常数，但双对数 ${\ell}_{\pm}^2/\bar{\ell}_{\pm}^2$ 中的虚部是依赖动量的。

式~(\ref{eq:matching})的一个特征行为是：该因子化是乘法型的~\cite{Ji:2020ect}，这表明硬胶子的贡献是定域的。其原因在于：准 TMDWFs 中夸克部分与反夸克部分之间的硬胶子交换是被幂次压低的——若存在这样的硬胶子，其两个连接点之间的空间间隔将远小于 $b_{\perp}$，相比典型硬模式的贡献而被幂次压低。
因此，在领头幂次下准 TMDWFs 的因子化是乘法型的。这一特征示于图~\ref{fig:reduced_graph}，其中共线、软与硬子图表示相应的贡献。此外，由共线动量模式的洛伦兹不变组合产生的 $\zeta_z$ 将为硬子图提供自然的硬标度。关于 LaMET 中准 TMDPDFs 因子化的更详细说明见近期综述~ \cite{Ji:2019ewn}。


%%%%%%%%%%%%%%%%%%%%%%%%%%%%%%%%%%
\begin{figure}
\centering
\includegraphics[width=0.5\textwidth]{reduced_graph}
\caption{赝标介子准 TMDWFs 的领头幂次约化图。此处 $C, ~S$ 表示红外结构中的共线与软部分，而 $H$ 表示硬贡献。由于夸克与反夸克之间的硬胶子交换被幂次压低，各硬部分彼此不相连，因此准波函数振幅的因子化是乘法型的。}
\label{fig:reduced_graph}
\end{figure}
%%%%%%%%%%%%%%%%%%%%%%%%%%%%%%%%%%




\subsection{由准 TMDWFs 确定 Collins-Soper 核}\label{sec:extractingCSkernel}


由式~(\ref{eq:matching})可见，准 TMDWFs 中的动量依赖提供了一条确定 CS 核的途径。可以将它写成与式~(\ref{eq:TMDsrapidityevolution})类似的形式~\cite{Ji:2021znw}：
\begin{align}
	2 \zeta_z \frac{d}{d \zeta_z} \ln \tilde \Psi^{\pm}\left(x, b_{\perp}, \mu, \zeta_z\right)=&K\left(b_{\perp}, \mu\right) +\frac{1}{2}\mathcal{G}^{\pm}\left(x^2\zeta_z,\mu\right) \nonumber\\
	+& \frac{1}{2}\mathcal{G}^{\pm}\left(\bar{x}^2\zeta_z,\mu\right)+\mathcal{O}\left(\frac{1}{\zeta_z}\right), \label{eq:quasiTMDsrapidityevolution}
\end{align}
其中 $K\left(b_{\perp}, \mu\right)$ 表示与式~(\ref{eq:TMDsrapidityevolution})中相同的核，且在大 $P^z$ 下不依赖硬标度 $\zeta_z$。与 TMDWFs 不同，准分布还包含硬贡献，其快度依赖由作为硬标度 $\zeta_z$ 函数的微扰项 $\mathcal{G}^{\pm}$ 表示。
由准 TMDWFs 的 $\zeta_z$ 依赖可以看出：当 $P^z\to\infty$ 时，$P^z$ 中的大对数被部分吸收进 $K\left(b_{\perp}, \mu\right)$，其余部分则并入微扰匹配 kernel。因此，要描述准 TMDWFs 对 $\zeta_z$ 的依赖，匹配 kernel $H^{\pm}$ 与 CS 核 $K\left(b_{\perp}, \mu\right)$ 的指数两者都需要。

为了显式地提取 CS 核 $K\left(b_{\perp}, \mu\right)$，可以利用式~(\ref{eq:matching})取两个不同的大动量 $P_1^z\neq P_2^z\gg1/b_{\perp}$ 但相同的标度 $\zeta_z$。取这两个量之比得到
\begin{align}
	\frac{\tilde \Psi^{\pm}(x, b_{\perp}, \mu,P_1^z)}{\tilde \Psi^{\pm}(x, b_{\perp}, \mu, P_2^z)}=\frac{H^{\pm}\left(xP_1^z,\mu \right)}{H^{\pm}\left(xP_2^z,\mu \right)}\exp\left[K(b_{\perp},\mu) \ln\frac{P_1^z}{P_2^z}\right], \label{eq:ratioofquasiWF}
\end{align}
其中约化软函数 $S_r(b_{\perp},\mu)$ 与 TMDWFs $\Psi^{\pm}(x,b_{\perp},\mu,\zeta)$ 已在该比值中相消。于是 CS 核 $K\left(b_{\perp}, \mu\right)$ 可以通过下式提取
\begin{align}
	K\left(b_{\perp}, \mu\right)=\frac{1}{\ln(P_1^z/P_2^z)} \ln\frac{H^{\pm}(xP_2^z,\mu)\tilde{\Psi}^{\pm}(x,b_{\perp},\mu,P_1^z)}{H^{\pm}(xP_1^z,\mu)\tilde{\Psi}^{\pm}(x,b_{\perp},\mu,P_2^z)}. \label{eq:extractingCSkernel}
\end{align}
注意提取的结果在领头幂次下形式上与 $x$ 及 $P_{1/2}^z$ 无关，并且 $\tilde{\Psi}^{+}$ 与 $\tilde{\Psi}^{-}$ 都可用于提取 $K\left(b_{\perp}, \mu\right)$。这是在该因子化方案的领头幂次下导出的，可能会被幂次修正破坏。因此，为了减小系统不确定度，我们取平均：
 \begin{align}
	K\left(b_{\perp}, \mu\right)=&\frac{1}{2\ln(P_1^z/P_2^z)} \left[\ln\frac{H^{+}(xP_2^z,\mu)\tilde{\Psi}^{+}(x,b_{\perp},\mu,P_1^z)}{H^{+}(xP_1^z,\mu)\tilde{\Psi}^{+}(x,b_{\perp},\mu,P_2^z)}\right.  \nonumber\\
	&+\left.\ln\frac{H^{-}(xP_2^z,\mu)\tilde{\Psi}^{-}(x,b_{\perp},\mu,P_1^z)}{H^{-}(xP_1^z,\mu)\tilde{\Psi}^{-}(x,b_{\perp},\mu,P_2^z)}\right].
\end{align}
具体细节将在第~\ref{subsec:csresults}节讨论。
